%% Đề 19 -- de-thi-chon-hsg-tinh-toan-9-nam-hoc-2019-2020-so-gddt-lang-son.pdf
\deheader{ĐỀ SỐ 19}%
  {Kỳ thi Chọn Học sinh giỏi cấp tỉnh lớp 9 — Năm học 2019--2020}%
  {Sở GD\&ĐT Lạng Sơn | Chỉ phần Đại số}
\addcontentsline{toc}{section}{Đề 19 — Kỳ thi Chọn HSG cấp tỉnh — Năm học 2019--2020}

\begin{de_cau}{Câu 1}{4,0 điểm}
Cho biểu thức:
\[A = \left(\dfrac{2}{\sqrt{ab}} + \dfrac{1}{a} + \dfrac{1}{b}\right) : \dfrac{\sqrt{a^3} + b\sqrt{a} + a\sqrt{b} + \sqrt{b^3}}{\sqrt{a^3b} + \sqrt{ab^3}},\quad \text{với } a > 0,\; b > 0.\]
\begin{enumerate}[label=\alph*)]
\item Rút gọn biểu thức $A$.
\item Biết $ab = 81$, tìm $a, b$ để $A$ đạt giá trị nhỏ nhất. Tìm giá trị nhỏ nhất đó.
\end{enumerate}
\end{de_cau}

\begin{loi_giai_cau}
\textbf{a)} Ta có:
\[\frac{2}{\sqrt{ab}} + \frac{1}{a} + \frac{1}{b} = \frac{2\sqrt{ab} + b + a}{ab} = \frac{(\sqrt{a}+\sqrt{b})^2}{ab}.\]
Biến đổi tử số của biểu thức chia:
\[\sqrt{a^3} + b\sqrt{a} + a\sqrt{b} + \sqrt{b^3} = \sqrt{a}(a+b) + \sqrt{b}(a+b) = (\sqrt{a}+\sqrt{b})(a+b).\]
Biến đổi mẫu số của biểu thức chia:
\[\sqrt{a^3b} + \sqrt{ab^3} = a\sqrt{ab} + b\sqrt{ab} = (a+b)\sqrt{ab}.\]
Do đó:
\[\frac{\sqrt{a^3} + b\sqrt{a} + a\sqrt{b} + \sqrt{b^3}}{\sqrt{a^3b} + \sqrt{ab^3}} = \frac{(\sqrt{a}+\sqrt{b})(a+b)}{(a+b)\sqrt{ab}} = \frac{\sqrt{a}+\sqrt{b}}{\sqrt{ab}}.\]
Suy ra:
\[A = \frac{(\sqrt{a}+\sqrt{b})^2}{ab} : \frac{\sqrt{a}+\sqrt{b}}{\sqrt{ab}} = \frac{(\sqrt{a}+\sqrt{b})^2}{ab} \cdot \frac{\sqrt{ab}}{\sqrt{a}+\sqrt{b}} = \boxed{\frac{\sqrt{a}+\sqrt{b}}{\sqrt{ab}} = \frac{1}{\sqrt{a}} + \frac{1}{\sqrt{b}}}.\]

\textbf{b)} Áp dụng bất đẳng thức AM-GM cho hai số dương $\dfrac{1}{\sqrt{a}}$ và $\dfrac{1}{\sqrt{b}}$:
\[A = \frac{1}{\sqrt{a}} + \frac{1}{\sqrt{b}} \ge 2\sqrt{\frac{1}{\sqrt{a}\cdot\sqrt{b}}} = \frac{2}{\sqrt[4]{ab}}.\]
Với $ab = 81$, ta có $\sqrt[4]{ab} = \sqrt[4]{81} = 3$, do đó:
\[A \ge \frac{2}{3}.\]
Dấu ``='' xảy ra khi và chỉ khi:
\[\frac{1}{\sqrt{a}} = \frac{1}{\sqrt{b}} \iff a = b.\]
Mà $ab = 81$ nên $a^2 = 81 \implies a = b = 9$ (vì $a, b > 0$).
Vậy $\min A = \boxed{\dfrac{2}{3}}$, đạt được khi $\boxed{a = b = 9}$.
\end{loi_giai_cau}

\begin{de_cau}{Câu 2}{4,0 điểm}
Cho phương trình: $x^2 + mx + m - 3 = 0$, với $m$ là tham số.
\begin{enumerate}[label=\alph*)]
\item Chứng minh rằng phương trình luôn có hai nghiệm phân biệt với mọi $m$.
\item Tìm giá trị của tham số $m$ để hai nghiệm $x_1, x_2$ của phương trình thỏa mãn:
\[\dfrac{1}{x_1^2} + \dfrac{1}{x_2^2} = \dfrac{2}{3}.\]
\end{enumerate}
\end{de_cau}

\begin{loi_giai_cau}
\textbf{a)} Phương trình có:
\[\Delta = m^2 - 4(m-3) = m^2 - 4m + 12 = (m-2)^2 + 8.\]
Vì $(m-2)^2 \ge 0$ với mọi $m$, nên $\Delta \ge 8 > 0$ với mọi $m$.
Vậy phương trình luôn có hai nghiệm phân biệt $x_1, x_2$ với mọi $m$.

\textbf{b)} Để biểu thức xác định thì $x_1, x_2 \neq 0 \iff m - 3 \neq 0 \iff m \neq 3$.
Theo định lý Vi-ét:
\[\begin{cases} x_1 + x_2 = -m \\ x_1 x_2 = m - 3 \end{cases}\]
Ta có:
\[\frac{1}{x_1^2} + \frac{1}{x_2^2} = \frac{x_1^2 + x_2^2}{x_1^2 x_2^2} = \frac{(x_1+x_2)^2 - 2x_1 x_2}{(x_1 x_2)^2} = \frac{(-m)^2 - 2(m-3)}{(m-3)^2} = \frac{m^2 - 2m + 6}{(m-3)^2}.\]
Theo đề bài:
\[\frac{m^2 - 2m + 6}{(m-3)^2} = \frac{2}{3} \iff 3(m^2 - 2m + 6) = 2(m^2 - 6m + 9).\]
\[\iff 3m^2 - 6m + 18 = 2m^2 - 12m + 18 \iff m^2 + 6m = 0 \iff m(m+6) = 0.\]
Suy ra $m = 0$ hoặc $m = -6$ (cả hai đều thỏa mãn $m \neq 3$).
Vậy $\boxed{m \in \{0; -6\}}$.
\end{loi_giai_cau}

\begin{de_cau}{Câu 3}{4,0 điểm}
\begin{enumerate}[label=\alph*)]
\item Cho $a, b, c, d$ là các số thực dương thỏa mãn $a+b+c+d = 4$. Chứng minh rằng:
\[\dfrac{1}{a^2+1} + \dfrac{1}{b^2+1} + \dfrac{1}{c^2+1} + \dfrac{1}{d^2+1} \ge 2.\]
\item Cho hình chữ nhật có độ dài hai cạnh là $2$ và $4$. Đặt vào bên trong hình chữ nhật đó $17$ điểm phân biệt, bất kỳ. Chứng minh rằng bao giờ cũng tìm được ít nhất ba điểm trong số 17 điểm đó tạo thành một tam giác có diện tích không vượt quá $\dfrac{1}{2}$.
\end{enumerate}
\end{de_cau}

\begin{loi_giai_cau}
\textbf{a)} Sử dụng kỹ thuật Cauchy ngược dấu:
\[\frac{1}{a^2+1} = 1 - \frac{a^2}{a^2+1} \ge 1 - \frac{a^2}{2a} = 1 - \frac{a}{2}.\]
Thiết lập tương tự cho $b, c, d$:
\[\frac{1}{b^2+1} \ge 1 - \frac{b}{2};\quad \frac{1}{c^2+1} \ge 1 - \frac{c}{2};\quad \frac{1}{d^2+1} \ge 1 - \frac{d}{2}.\]
Cộng vế theo vế cả 4 bất đẳng thức:
\[\sum \frac{1}{a^2+1} \ge 4 - \frac{a+b+c+d}{2} = 4 - \frac{4}{2} = \boxed{2}.\]
Dấu ``='' xảy ra khi và chỉ khi $a = b = c = d = 1$. Điều phải chứng minh.

\textbf{b)} Chia hình chữ nhật $2 \times 4$ thành 8 hình vuông nhỏ cạnh $1 \times 1$.
Diện tích mỗi hình vuông nhỏ là $1$.
Theo nguyên lý Dirichlet, khi đặt 17 điểm vào 8 hình vuông nhỏ:
\[\left\lfloor \frac{17-1}{8} \right\rfloor + 1 = 2 + 1 = 3.\]
Do đó, tồn tại ít nhất một hình vuông nhỏ chứa ít nhất 3 điểm.
Gọi 3 điểm đó là $M, N, P$. Tam giác $MNP$ nằm hoàn toàn trong hình vuông cạnh 1.
Diện tích của một tam giác nằm trong hình vuông cạnh 1 luôn không vượt quá nửa diện tích hình vuông đó:
\[S_{MNP} \le \frac{1}{2} S_{\text{hình vuông}} = \frac{1}{2} \cdot 1 = \boxed{\frac{1}{2}}.\]
Điều phải chứng minh.
\end{loi_giai_cau}

\begin{de_cau}{Câu 5}{2,0 điểm}
Tìm tất cả các cặp số nguyên $(x; y)$ thỏa mãn phương trình:
\[x^2 - 2xy + 2y^2 - 2x + 6y + 5 = 0.\]
\end{de_cau}

\begin{loi_giai_cau}
Biến đổi phương trình về dạng tổng các bình phương:
\[(x^2 - 2xy + y^2) - 2(x - y) + 1 + (y^2 + 4y + 4) = 0\]
\[\iff (x - y)^2 - 2(x - y) + 1 + (y + 2)^2 = 0\]
\[\iff (x - y - 1)^2 + (y + 2)^2 = 0.\]
Vì $(x - y - 1)^2 \ge 0$ và $(y + 2)^2 \ge 0$ với mọi $x, y \in \mathbb{R}$, nên tổng của chúng bằng 0 khi và chỉ khi:
\[\begin{cases} y + 2 = 0 \\ x - y - 1 = 0 \end{cases} \iff \begin{cases} y = -2 \\ x = y + 1 = -2 + 1 = -1 \end{cases}\]
Vậy phương trình có nghiệm nguyên duy nhất là $\boxed{(x; y) = (-1; -2)}$.
\end{loi_giai_cau}
